\section{Wprowadzenie}

Miód jest produktem, którego jakość ocenia się nie tylko za pomocą analiz
instrumentalnych, lecz także zmysłami: wzrokiem, węchem i~smakiem. Ocena
organoleptyczna jest tania i~szybko wykonalna, a~jej wynik ma znaczenie
praktyczne, ponieważ może zdecydować o~dopuszczeniu partii do obrotu.
Jednocześnie jest ona najtrudniejsza do ujednolicenia, bo opiera się na
ocenie człowieka, a~nie na odczycie przyrządu. Ramy prawne tej oceny
wyznaczają przepisy unijne, krajowe oraz normy międzynarodowe.

Jakość handlowa miodu dopuszczonego do obrotu w~Unii Europejskiej jest
określona dyrektywą 2001/110/WE \cite{Dyrektywa2001110}. Dyrektywa jest
implementowana w~prawie krajów członkowskich, a~poza Unią punktem odniesienia
jest norma Codex Alimentarius CXS 12-1981 \cite{CodexHoney}. Oba dokumenty,
obok parametrów oznaczanych instrumentalnie, zawierają wymagania
organoleptyczne. Miód nie może mieć obcego smaku ani zapachu, a~nazwę
odmianową można stosować tylko wtedy, gdy produkt wykazuje cechy
organoleptyczne, fizykochemiczne i~mikroskopowe właściwe dla deklarowanego
źródła.

Przepisy polskie stawiają wymaganie szersze. Rozporządzenie w~sprawie
jakości handlowej miodu \cite{RozporzadzenieMRiRW2003} uznaje za zgodny miód,
który ,,nie posiada zapachu i~smaku nietypowego dla danej odmiany miodu''.
Smak obcy pochodzi spoza miodu, na przykład z~substancji wchłoniętych podczas
przetwarzania i~przechowywania. Stwierdzenie smaku nietypowego wymaga
natomiast porównania z~wzorcem odmiany, którego żaden z~wymienionych aktów
nie definiuje. Charakterystyki opisowe głównych europejskich miodów
odmianowych opracowano w~literaturze \cite{PersanoOddo2004}, a~metodykę ich
oceny sensorycznej w~ramach International Honey Commission
\cite{Piana2004}, lecz nie mają one statusu wzorca prawnego. Jednocześnie
art.~5 dyrektywy zakazuje państwom członkowskim przyjmowania przepisów
krajowych dotyczących miodu, których dyrektywa nie przewiduje.

Metodę oceny organoleptycznej określa w~Polsce rozporządzenie w~sprawie metod
analiz miodu \cite{RozporzadzenieMRiRW2009}. Ogranicza się ono do wskazania,
że barwę, konsystencję, smak i~zapach bada się narządami wzroku, węchu
i~smaku, a~wyniki zestawia w~tabeli. Nie określa liczby ani kwalifikacji
oceniających, materiału odniesienia ani kryterium uznania smaku za
nietypowy. Dyrektywa (UE) 2024/1438, stosowana od 14~czerwca 2026~r.
\cite{Dyrektywa20241438}, przewiduje przyjęcie zharmonizowanych metod
wykrywania zafałszowań miodu, lecz nie metody oceny organoleptycznej.

Sposób przeprowadzenia oceny i~interpretacji jej wyniku pozostaje więc
w~gestii laboratorium. W~rezultacie werdykt organoleptyczny może samodzielnie
przesądzać o~niezgodności partii, nawet gdy pozostałe oznaczenia nie budzą
zastrzeżeń.

Od kryterium o~skutkach prawno-administracyjnych oczekuje się
obiektywności, powtarzalności i~proporcjonalności. W~kontroli urzędowej
wymóg ten ma podstawę prawną: metody badań laboratoryjnych muszą być zgodne
z~przepisami Unii, a~w~ich braku z~protokołami uznanymi międzynarodowo,
z~przepisami krajowymi lub zwalidowane w~badaniach między- lub
wewnątrzlaboratoryjnych \cite[art.~34]{Rozporzadzenie2017625}.
W~skoordynowanym planie kontroli miodu Komisji Europejskiej z~2015~r. test
sensoryczny włączono do pierwszego etapu charakterystyki próbek, obok
przewodności elektrycznej, aktywności diastazy i~analizy pyłkowej, jako
badanie wskazujące kierunek dalszych, bardziej specyficznych oznaczeń
\cite{KE2015PlanKontroli}. Mierzalność cechy nie
przesądza jeszcze o~jej przydatności jako kryterium decyzyjnego: wykazano to
dla mikrobiologicznych parametrów miodu, gdzie sama liczba drożdży
osmofilnych, bez oznak aktywnej fermentacji, nie stanowi wystarczającej
podstawy do zakwestionowania produktu \cite{Gajda2026}. Analogiczne pytanie
można postawić wobec oceny typowości smaku i~zapachu.

Literatura psychofizyczna wskazuje na dwa niezależne źródła zmienności takiej
oceny. Pierwszym jest międzyosobnicza, częściowo dziedziczna zmienność
percepcji smaków podstawowych: próg rozpoznania kwaśności jest dziedziczny
w~około 53\% \cite{wise2007twin}, a~polimorfizmy genu \textit{GNAT3}
wyjaśniają około 13\% wariancji percepcji słodyczy sacharozy
\cite{fushan2010sweet}. Ta sama próbka wywołuje więc u~różnych osób wrażenia
o~realnie różnej intensywności, a~doświadczeń zmysłowych nie da się
bezpośrednio porównać \cite{bartoshuk2000comparing}. Drugim źródłem jest
nieporównywalność skal opisowych: etykiety słowne (np.\ ,,bardzo silny'',
,,typowy'') oznaczają dla różnych osób różne intensywności bezwzględne,
co czyni skale kategorialne niewłaściwym narzędziem porównań między osobami
i~grupami \cite{bartoshuk2003labeled}.

Celem badania było empiryczne sprawdzenie, czy ocena typowości smaku i~zapachu
miodu jest na tyle powtarzalna, by mogła stanowić samodzielne kryterium
oceny jakości handlowej.
W~tym celu porównano wyniki ślepego testu konsumenckiego z~ocenami 
organoleptycznymi tych samych próbek, przeprowadzonymi niezależnie 
przez laboratoria. Analiza koncentrowała się na zdolności panelu konsumenckiego do różnicowania próbek pod względem typowości smaku, zgodności ocen między laboratoriami oraz stopniu, w~jakim konsumenci i~laboratoria wskazywali te same próbki jako nietypowe.
